% !TeX root = ../../Book3.tex
% 纲要标题：### 12.3 Fitting 理想、纤维生成元数与局部自由
% 纲要位置：计划纲要.md，第 1860 行。

\section{Fitting 理想、纤维生成元数与局部自由}
\label{b3:sec:1203}

% TODO: 按以下纲要要点撰写本节。

% 纲要内容（仅作写作计划，尚非教材正文）：
%
% * 从有限呈示矩阵的子式定义 Fitting 理想，证明其与呈示选择无关，并与任意基变换相容
% * 由剩余域上矩阵的秩解释 \(\dim_{\kappa(\mathfrak p)}(M\otimes_R\kappa(\mathfrak p))\) 的闭条件与上半连续性
% * 在有限呈示假设下建立局部自由与常秩的适当 Fitting 判据；区别纤维维数恒定与无附加假设的平坦性断言
% * 连接第二卷第4.2节的子式理想和本章泛自由性；一般环上保留理想不变量，不把它们误写成必然存在的 Smith 对角形
% * 对有限呈示模，局部自由秩为 r 的判据写成局部化后的 \(\operatorname{Fitt}_{r-1}(M)=0\)、\(\operatorname{Fitt}_{r}(M)=R\)（约定 \(\operatorname{Fitt}_{-1}=0\)）；用幂零基底上的例子区分这些理想等式与仅在剩余域上看到的维数恒定。
%

待撰写。
